\chapter{Segment naziemny}\label{ch:naziemny}

\begin{sekcja}{Stanowisko startowe}
Wybierz otwarty teren bez drzew i linii energetycznych, np. boisko. Uzyskaj zgodę właściciela terenu. Powłoka jest delikatna: gałązki, korzenie i kamienie mogą ją przebić.

Rysunek~\ref{fig:stanowisko} pokazuje warstwy stanowiska. Powłoka leży \textbf{między dwiema foliami}. Tankowanie odbywa się między foliami. Folia górna przytrzymuje powłokę, aby nie odleciała przed startem.

\begin{center}\resizebox{\linewidth}{!}{\rysStanowisko}\end{center}
\captionof{figure}{Przekrój stanowiska startowego (grubość warstw przesadzona).}\label{fig:stanowisko}
\end{sekcja}

\begin{sekcja}{Warstwy stanowiska}
\begin{enumerate}
  \item \textbf{Plandeka} -- na sam dół. Chroni przed wilgocią z ziemi. Jest ważna przy starcie rano, gdy trawa jest mokra.
  \item \textbf{Koc} -- chroni powłokę przed przebiciem przez korzenie i odrosty. Gruby koc działa lepiej niż cienki.
  \item \textbf{Folia malarska dolna} (cienka) -- pod powłoką.
  \item \textbf{Powłoka} i rura do tankowania.
  \item \textbf{Folia malarska górna} (grubsza) -- na powłoce. Cienka folia na wierzchu może się porwać i powłoka może odlecieć w trakcie tankowania. Grubsza folia jest pewniejsza.
\end{enumerate}

\begin{uwaga}Dociśnij rogi plandeki, koca i folii czymś ciężkim, np. skrzynkami transportowymi. Porywy wiatru mogą porwać folię razem z powłoką.\end{uwaga}

\begin{wskazowka}Folię, plandekę i koce kupisz w markecie budowlanym. Kup oba rodzaje folii, aby mieć wybór. Kup wszystko na jeden rachunek razem z zamówieniem gazu -- upraszcza to rozliczenie.\end{wskazowka}
\end{sekcja}

\begin{sekcja}{Stanowisko na zdjęciach}
\dwafoto{2023_stanowisko_rozlozone}{Rozłożona plandeka i sprzęt przed tankowaniem.}
{2023_powloka_miedzy_foliami}{Powłoka na folii dolnej; pomocnicy przykrywają ją folią górną.}
\dwafoto{2022_stanowisko_zabezpieczone}{Rogi folii dociśnięte skrzynkami transportowymi.}
{2023_balon_po_tankowaniu}{Zatankowana powłoka trzymana folią górną do chwili startu.}
\end{sekcja}

\begin{sekcja}{Sprzęt do tankowania}
\begin{itemize}
  \item butla z helem (zwykle \SI{50}{l}, \SI{200}{bar}) i reduktor z dwoma manometrami,
  \item wąż i rura do tankowania (najlepiej z zaworem zwrotnym),
  \item obciążnik o masie równej sile nośnej szyjki -- dla ładunku do \SI{4}{kg} to kilka kilogramów: worek z piaskiem, plecak z obciążeniem albo odważniki (butelka z wodą jest za lekka),
  \item waga bagażowa (zakres do kilkunastu kilogramów) do ważenia obciążnika i zestawu,
  \item klucz do reduktora: klucz nastawny („francuski”) albo klucze calowe -- zależnie od butli,
  \item taśma klejąca (mocna, np. srebrna), sznurek, opaski zaciskowe,
  \item rękawiczki (bawełniane albo nitrylowe) dla wszystkich, którzy dotykają powłoki,
  \item termometry do pomiaru temperatury butli i powietrza (sekcja~\ref{sec:temp-butli}): pirometr, czarna matowa taśma izolacyjna, termometr pogodowy, opcjonalnie termometr kontaktowy,
  \item czujnik wodoru -- na wszelki wypadek, gdy gaz w butli może zawierać wodór albo gdy zespół używa wodoru.
\end{itemize}

\begin{uwaga}Pracuj przy powłoce tylko w rękawiczkach. Tłuszcz ze skóry osłabia lateks. Powłoka pęka wcześniej w miejscach, które ktoś dotknął gołą ręką.\end{uwaga}

\begin{uwaga}Butla z gazem pod ciśnieniem \SI{200}{bar} jest niebezpieczna~\cite[rozdz.~6]{nasahb}:
\begin{itemize}
  \item Transportuj ją zabezpieczoną przed przetoczeniem (rozdział~\ref{ch:planowanie}). Nie upuszczaj jej i nie uderzaj.
  \item Nie tocz i nie ciągnij butli. Na krótkie odległości przewoź ją na wózku.
  \item Nie dopuszczaj do nagrzania powyżej \SI{50}{\degreeCelsius}: nie zostawiaj butli w samochodzie na słońcu.
  \item Nie zdejmuj kołpaka z zaworu, dopóki butla nie leży stabilnie na stanowisku.
  \item Nie otwieraj zaworu bez zamontowanego reduktora.
  \item Przed misją obejrzyj wąż, reduktor i uszczelki. Nie zginaj węża na ostro -- zagięcie osłabia go niewidocznie.
\end{itemize}\end{uwaga}

\begin{uwaga}[Transport butli]Przewóz butli z gazem sprężonym podlega przepisom o przewozie towarów niebezpiecznych (ADR) i może wymagać uprawnień. Jeżeli procedury zakupowe na to pozwalają, kup albo wypożycz butlę u dostawcy blisko miejsca startu i odbierz ją na miejscu.\end{uwaga}
\end{sekcja}

\begin{sekcja}{Przygotowanie reduktora}
Reduktor ma dwa manometry:
\begin{itemize}
  \item \textbf{manometr wysokiego ciśnienia} (ciśnienie w butli, skala do ok. \SI{300}{bar}) -- z niego odczytujesz $p_1$ i $p_2$,
  \item \textbf{wskaźnik przepływu} (skala w l/min, skalowana dla argonu i CO$_2$) -- do ustawienia i utrzymania stałego przepływu.
\end{itemize}

\begin{enumerate}
  \item Połóż butlę stabilnie na plandece, w cieniu. Zabezpiecz ją przed przetoczeniem. Naklej na butlę kawałek czarnej taśmy do pomiaru temperatury (sekcja~\ref{sec:temp-butli}).
  \item Sprawdź gwint zaworu butli i reduktora. Usuń zanieczyszczenia.
  \item Przykręć reduktor kluczem. Nie dokręcaj na siłę.
  \item Zamknij pokrętło reduktora (pozycja \emph{OFF}). Otwórz powoli zawór butli.
  \item Odczytaj z \textbf{manometru wysokiego ciśnienia} ciśnienie początkowe $p_1$ (bez przepływu -- reduktor zamknięty). Zmierz temperaturę butli $T_1$. Zapisz obie wartości w formularzu F10.
  \item Sprawdź szczelność: zamknij zawór butli. Ciśnienie na manometrze wysokiego ciśnienia nie może spadać.
  \item Podłącz wąż i rurę do tankowania.
\end{enumerate}

\dwafoto{2022_regulator_montaz}{Montaż reduktora na butli -- w rękawiczkach, na plandece.}
{2022_regulator_manometry}{Manometry reduktora: górny -- ciśnienie w butli [bar], dolny -- przepływ [l/min], skalowany dla argonu i CO$_2$.}
\end{sekcja}

\begin{sekcja}{Ilość gazu z ciśnienia w butli}\label{sec:gaz-wzor}
Manometr wysokiego ciśnienia pokazuje ciśnienie w butli. Ze spadku ciśnienia obliczysz, ile gazu poszło do powłoki. Wzór wynika z równania stanu gazu z poprawką na ściśliwość helu~\cite{nist}, \cite[rozdz.~6]{nasahb}:
\[
  V \;=\; \frac{V_b\,T_0}{p_0}\left(\frac{p_1}{Z(p_1)\,T_1} - \frac{p_2}{Z(p_2)\,T_2}\right),
  \qquad Z(p) \approx 1 + 4{,}8\cdot10^{-4}\,p\ [\text{bar}]
\]
gdzie:
\begin{itemize}
  \item $V$ -- objętość helu wypuszczonego do powłoki, przy ciśnieniu atmosferycznym [m³],
  \item $V_b$ -- pojemność wodna butli (wybita na butli, np. \SI{50}{l} = \SI{0.05}{m^3}),
  \item $p_1$, $p_2$ -- ciśnienie \textbf{bezwzględne} przed i po tankowaniu: odczyt manometru wysokiego ciśnienia + \SI{1}{bar}; odczytuj tylko przy zamkniętym zaworze -- w czasie przepływu manometr pokazuje za mało,
  \item $p_0 = \SI{1.013}{bar}$ -- ciśnienie atmosferyczne,
  \item $Z$ -- współczynnik ściśliwości helu (dla \SI{200}{bar} $Z \approx 1{,}1$; bez poprawki wynik jest zawyżony o około \SI{10}{\percent}),
  \item $T_1$, $T_2$ -- temperatura gazu w butli [K] przy odczycie $p_1$ i $p_2$ (sekcja~\ref{sec:temp-butli}); $T_0 = \SI{288}{K}$. Błąd temperatury \SI{10}{\degreeCelsius} to błąd ilości gazu około \SI{3.5}{\percent}.
\end{itemize}

Siła nośna helu przy ziemi to około \SI{1.05}{kg} na każdy \si{m^3}. Siła nośna szyjki = $V \cdot 1{,}05$~kg $-$ masa powłoki.
\end{sekcja}

\begin{sekcja}{Tabela gazu}
\small
\begin{tabularx}{\linewidth}{@{}YYYYY@{}}
\toprule
\makecell{\textbf{Odczyt}\\\textbf{manometru [bar]}} & \makecell{\textbf{Spadek}\\\textbf{ciśnienia [bar]}} & \makecell{\textbf{Hel z butli}\\\textbf{10 l [m³]}} & \makecell{\textbf{Hel z butli}\\\textbf{50 l [m³]}} & \makecell{\textbf{Siła nośna}\\\textbf{(50 l) [kg]}}\\\midrule
\wiersze{figures/generated/tabela_gazu.tex}
\bottomrule
\end{tabularx}
\normalsize
\captionof{table}{Hel wypuszczony z butli napełnionej do \SI{200}{bar}, w temperaturze \SI{15}{\degreeCelsius}. Pełna butla \SI{50}{l} zawiera około \SI{9}{m^3} helu. Reguła: z butli \SI{50}{l} każde \SI{10}{bar} spadku to około 0,4--0,5~m³ helu.}\label{tab:gaz}

\begin{uwaga}Przy szybkim wypływie butla się ochładza i manometr pokazuje za niskie ciśnienie. Przed odczytem $p_2$ zamknij zawór i odczekaj kilka minut, aż temperatura się wyrówna.\end{uwaga}

Hel balonowy („balloon gas”) bywa mieszaniną helu z powietrzem. Ma mniejszą siłę nośną. Zamów hel techniczny o czystości podanej przez dostawcę i zapisz ją w karcie misji.
\end{sekcja}

\begin{sekcja}{Pomiar temperatury butli}\label{sec:temp-butli}
Termometr mierzy temperaturę \textbf{ścianki} butli, nie gazu. Ścianka i gaz mają tę samą temperaturę tylko w równowadze: butla stoi długo w stałej temperaturze, w cieniu, a gaz nie płynie~\cite[rozdz.~6]{nasahb}. W czasie wypływu gaz w butli się ochładza, a słońce nagrzewa ściankę -- pomiar jest wtedy fałszywy.

\textbf{Przygotowanie:} od przyjazdu trzymaj butlę w cieniu (albo przykrytą jasnym materiałem). Naklej na butlę, na wysokości połowy, kawałek \textbf{czarnej matowej taśmy izolacyjnej PVC} (około 5$\times$\SI{5}{cm}). Taśma ma znaną emisyjność $\varepsilon \approx 0{,}95$. Naklej ją kilka minut przed pomiarem, aby przyjęła temperaturę butli.

\textbf{Pirometr} (termometr na podczerwień) -- metoda podstawowa:
\begin{enumerate}
  \item Ustaw w pirometrze emisyjność $\varepsilon = 0{,}95$ (wiele tanich pirometrów ma ją ustawioną na stałe).
  \item Celuj w taśmę, nie w gołą butlę. Goły lub błyszczący metal ma małą i nieznaną emisyjność: odbija promieniowanie otoczenia i pirometr pokazuje fałszywą wartość.
  \item Podejdź tak blisko, aby plamka pomiaru była mniejsza niż taśma. Przy pirometrze 12:1 z odległości \SI{12}{cm} plamka ma około \SI{1}{cm}.
  \item Zmierz trzy razy i zapisz średnią.
\end{enumerate}

\textbf{Termometr kontaktowy} (sonda powierzchniowa, np. termopara typu K z multimetrem) -- kontrola: dociśnij sondę do butli, przykryj ją kawałkiem pianki albo rękawiczką i odczekaj, aż odczyt się ustali.

\textbf{Termometr pogodowy} -- temperatura powietrza w cieniu. Gdy butla stała kilka godzin na zewnątrz w cieniu, jej temperatura jest bliska temperaturze powietrza. Duża różnica oznacza błąd pomiaru albo butlę nagrzaną w samochodzie.

\textbf{Kiedy mierzyć:} $T_1$ -- przed otwarciem reduktora, razem z $p_1$. $T_2$ -- po zamknięciu zaworu i odczekaniu kilku minut, razem z $p_2$. Zapisz obie wartości i temperaturę powietrza w formularzu F10.
\end{sekcja}

\begin{sekcja}{Hel czy wodór}\label{sec:hel-wodor}
Zespół wybiera gaz na etapie planowania misji (sekcja~\ref{sec:powloka}). Platforma Stardust NEXT używa \textbf{helu}. Hel drożeje i bywa trudno dostępny, dlatego część zespołów balonowych przechodzi na wodór~\cite[rozdz.~2]{nasahb}. Wodór daje około \SI{10}{\percent} więcej siły nośnej, ale jest palny i wymaga dodatkowego szkolenia.

\begin{uwaga}Nowym zespołom nie zalecamy wodoru. Jeśli zespół mimo to wybierze wodór, stosuj zasady z rozdziału~6 podręcznika NASA~\cite{nasahb}:
\begin{itemize}
  \item żadnego ognia, iskier ani testów elektroniki przy butli i przy powłoce; silniki samochodów wyłączone,
  \item uziem butlę i sprzęt, zanim wypuścisz gaz,
  \item odzież z bawełny, nie z tworzyw sztucznych,
  \item przed tankowaniem wyznacz strefę zagrożenia: owal wzdłuż wiatru, balon przy jego nawietrznym końcu. Uciekaj pod wiatr, gdy stoisz po nawietrznej; z innych miejsc -- w poprzek wiatru. Nigdy nie przechodź przez oś owalu,
  \item płomień wodoru jest niewidoczny. Ktoś z zewnątrz strefy obserwuje stanowisko i ostrzega,
  \item resztki powłoki po locie mogą zawierać kieszenie gazu -- nie zostawiaj ich bez nadzoru.
\end{itemize}
Wzór i tabela gazu w tym rozdziale dotyczą tylko helu.\end{uwaga}
\end{sekcja}

\begin{sekcja}{Ilość gazu z przepływu i czasu}
Operator butli może szacować ilość gazu na bieżąco: ustaw \textbf{stały przepływ} na wskaźniku przepływu i mierz czas.

Wskaźnik przepływu jest skalowany dla argonu. Hel jest około 10 razy lżejszy, więc rzeczywisty przepływ helu jest około \textbf{3 razy większy} niż wskazanie ($\sqrt{40/4} \approx 3{,}2$).

\textbf{Reguła inżynierska:}
\[
  t\,[\text{min}] \;\approx\; \frac{310 \cdot V\,[\text{m}^3]}{Q_\text{wsk}\,[\text{l/min}]}
\]
Przykład: przy wskazaniu \SI{30}{l/min} tankowanie trwa około \textbf{10 minut na każdy m³} helu.

\renewcommand{\arraystretch}{1.25}
\small
\begin{tabularx}{\linewidth}{@{}YYYY@{}}
\toprule
\textbf{Wskazanie [l/min]} & \textbf{Hel rzeczywiście [l/min]} & \textbf{Czas na 1 m³ [min]} & \textbf{Czas na 3 m³ [min]}\\\midrule
10 & ok. 32 & 31 & 93\\
20 & ok. 63 & 16 & 47\\
30 & ok. 95 & 10 & 31\\
\bottomrule
\end{tabularx}
\normalsize

To zgrubne oszacowanie: wskazanie zależy od reduktora i od ciśnienia za nim. Co kilka minut porównuj je ze spadkiem ciśnienia w butli (tabela~\ref{tab:gaz}). O końcu tankowania decyduje metoda z obciążnikiem.
\end{sekcja}

\begin{sekcja}{Metoda z obciążnikiem}
Metoda z obciążnikiem to podstawowa metoda tankowania. Wzory służą do kontroli i do oceny, ile gazu zostało w butli.
\begin{enumerate}
  \item Odczytaj z kalkulatora siłę nośną szyjki (\emph{neck lift}).
  \item Przygotuj obciążnik o masie równej \textbf{sile nośnej szyjki z kalkulatora} (worek z piaskiem, plecak, odważniki). Uwzględnij masę rury do tankowania i taśmy, jeśli wiszą razem z obciążnikiem. Zważ obciążnik wagą bagażową.
  \item Przywiąż obciążnik do szyjki albo do rury do tankowania.
  \item Tankuj, aż balon \textbf{uniesie obciążnik} (balon w równowadze albo lekko się wznosi).
\end{enumerate}

\begin{uwaga}Masa obciążnika to siła nośna szyjki z kalkulatora, \textbf{nie} masa gondoli ze zgłoszenia do PAŻP. Siła nośna szyjki = masa zestawu + nadwyżka siły nośnej, która daje prędkość wznoszenia. Balon zatankowany „na masę gondoli” prawie nie ma nadwyżki: wyrwie się z rąk, a potem będzie się wlókł nisko -- po budynkach, liniach energetycznych i drzewach. To skrajnie niebezpieczne.\end{uwaga}

\begin{center}\resizebox{\linewidth}{!}{\rysTankowanie}\end{center}
\captionof{figure}{Tankowanie metodą z obciążnikiem.}
\end{sekcja}

\begin{sekcja}{Procedura tankowania}
\begin{enumerate}
  \item Wszyscy przy powłoce zakładają rękawiczki.
  \item Wyjmij powłokę z opakowania na folii dolnej. Rozłóż ją bez naciągania.
  \item Włóż rurę do tankowania w szyjkę powłoki. Owiń szyjkę dookoła rury i przytrzymaj ręką.
  \item Wstępnie uszczelnij szyjkę na rurze taśmą -- \textbf{delikatnie}, tak aby dało się ją zdjąć.
  \item Przykryj powłokę folią górną. Pomocnicy trzymają brzegi folii.
  \item Otwórz zawór butli i powoli otwieraj reduktor. Ustaw stały przepływ. Zapisz godzinę.
  \item Obserwuj manometr wysokiego ciśnienia. Zapisuj odczyty (formularz F10). Porównuj z tabelą~\ref{tab:gaz} i z czasem tankowania.
  \item Gdy balon uniesie obciążnik, zamknij zawór butli. Sprawdź siłę nośną jeszcze raz po minucie.
  \item Po kilku minutach odczytaj $p_2$ i zmierz temperaturę butli $T_2$ (formularz F10).
  \item Zamknij powłokę (sekcja~\ref{sec:szyjka}).
\end{enumerate}

\dwafoto{2023_tankowanie_rura}{Tankowanie: szyjka owinięta dookoła rury i przytrzymana ręką. Pętla olinowania z krętlikiem gotowa. Taśma w gotowości na plandece.}
{2020_tankowanie_torun}{Tankowanie pod folią (20.07.2020).}
\end{sekcja}

\begin{sekcja}{Zamknięcie powłoki}\label{sec:szyjka}
Przygotuj wcześniej pętlę olinowania z krętlikiem. Pętla jest gotowa do założenia zaraz po tankowaniu.
\begin{enumerate}
  \item Zdejmij taśmę wstępną. Wysuń rurę z szyjki, trzymając szyjkę zaciśniętą w dłoni.
  \item Skręć szyjkę kilka razy.
  \item Przewiąż skręconą szyjkę sznurkiem. Zrób kilka supłów.
  \item Przełóż przez szyjkę przygotowaną pętlę olinowania (z krętlikiem), \textbf{przed} zgięciem szyjki.
  \item Zegnij szyjkę w pół \textbf{w górę} (w stronę powłoki). Pętla zostaje w zgięciu.
  \item Zaciśnij zgiętą szyjkę kilkoma opaskami zaciskowymi.
  \item Owiń całość kilkoma warstwami taśmy.
  \item Pociągnij za pętlę: pętla nie może się zsunąć z szyjki.
\end{enumerate}

Zamknięcie musi być \textbf{szczelne}. Szczelności nie da się łatwo sprawdzić na stanowisku -- zdaj się na intuicję inżynierską: wiąż mocno, nie żałuj supłów, opasek i taśmy. Nieszczelna szyjka to wolniejsze wznoszenie, niższy pułap albo lot bez pęknięcia.

\begin{uwaga}Nie przyklejaj taśmy do powłoki. Taśma odklejana od powłoki może ją przerwać. Taśma ma dotykać tylko szyjki.\end{uwaga}
\end{sekcja}

\begin{sekcja}{Stacje odbiorcze}
Używaj co najmniej dwóch stacji odbiorczych: stacjonarnej i mobilnej. Nie polegaj tylko na sieci innych nasłuchowców. Schemat toru radiowego: rozdział~\ref{ch:elektronika}.

\textbf{Stacja stacjonarna:} antena kierunkowa albo dookólna na wysokim maszcie, odbiornik SDR i komputer z internetem. Stacja wysyła odebrane pakiety do SondeHub i APRS-IS.

\textbf{Stacja mobilna}, wariant 1: radiotelefon z dekoderem APRS (np. Yaesu FT5DE) i antena samochodowa VHF/UHF na magnes (np. Nagoya UT-72).

\textbf{Stacja mobilna}, wariant 2: odbiornik RTL-SDR, przedwzmacniacz niskoszumny (LNA), filtry pasmowe, laptop z oprogramowaniem dekodującym.

\begin{wskazowka}[APRS bez dekodera]Tani radiotelefon bez dekodera APRS (np. Baofeng) też wystarczy. Połącz wyjście audio radia z telefonem albo laptopem (kabel audio albo adapter Bluetooth). Dekoduj APRS aplikacją APRSdroid (Android) albo programem Direwolf (laptop).\end{wskazowka}

\foto[0.6\linewidth]{2022_poscig_samochod}{Stacja mobilna w samochodzie pościgowym: laptop z mapą śledzenia, radiotelefon.}
\end{sekcja}

\begin{sekcja}{Konfiguracja stacji naziemnej}
Skonfiguruj stacje i sprawdź je z działającym trackerem najpóźniej 2 dni przed startem (formularz F7).

\renewcommand{\arraystretch}{1.3}
\small
\begin{tabularx}{\linewidth}{@{}P{26mm}P{36mm}L@{}}
\toprule
\textbf{Sygnał} & \textbf{Oprogramowanie} & \textbf{Konfiguracja}\\\midrule
RTTY \SI{434}{MHz} & horus-gui (Project Horus) albo dl-fldigi & Ustaw SDR na częstotliwość trackera, tryb RTTY (prędkość i bity jak w trackerze). Włącz wysyłanie do SondeHub Amateur ze swoim znakiem i pozycją stacji.\\
APRS \SI{144.8}{MHz} & Direwolf (laptop), APRSdroid (Android) & Direwolf jako iGate: wpisz znak i hasło APRS-IS, aby przekazywać pakiety do aprs.fi.\\
Horus Binary (opcjonalnie) & horus-gui & Format Horus Binary trafia bezpośrednio do SondeHub Amateur.\\
\bottomrule
\end{tabularx}
\normalsize

\textbf{Stacja SatNOGS.} Stacja sieci SatNOGS (antena, SDR, komputer) nadaje się na stację stacjonarną. Na czas misji zatrzymaj klienta SatNOGS i uruchom na tym samym SDR oprogramowanie z tabeli. Po misji przywróć klienta SatNOGS.

\textbf{Test:} włącz tracker w odległości co najmniej \SI{100}{m} od stacji. Sprawdź, czy pakiety ze stacji pojawiają się w SondeHub i na aprs.fi pod znakiem stacji.
\end{sekcja}

\begin{sekcja}{Łączność w zespole i poszukiwania}
Radiotelefony to podstawa łączności w czasie pościgu (między samochodami) i w terenie, przy szukaniu gondoli. Telefon komórkowy nie zawsze ma zasięg.

\begin{itemize}
  \item Używaj radiotelefonów \textbf{PMR446} -- każdy może z nich korzystać bez licencji.
  \item Ustal przed startem kanał główny i zapasowy. Zapisz je w karcie misji (F0).
  \item Każdy samochód i każda grupa poszukiwawcza ma co najmniej jeden radiotelefon.
\end{itemize}

\textbf{Przeszukiwanie terenu}, gdy ostatnia pozycja jest niedokładna:
\begin{enumerate}
  \item Wyznacz obszar wokół ostatniej pozycji (z trackera, SPOT, znaczników Bluetooth) z uwzględnieniem kierunku wiatru.
  \item Ustaw kilka osób w linii, w równych odstępach (w lesie mniejszych, na polu większych).
  \item Idźcie równolegle, w jednym kierunku. Utrzymujcie odstępy i kontakt wzrokowy albo radiowy.
  \item Na końcu pasa przesuńcie linię o jej szerokość i wróćcie. Zaznaczajcie przeszukane pasy na mapie.
  \item Szukajcie też w górze: gondola może zawisnąć na drzewie.
  \item Co kilkadziesiąt metrów zatrzymajcie się i nasłuchujcie brzęczyka gondoli w ciszy (rozdział~\ref{ch:elektronika}).
  \item Uważajcie na teren prywatny, zwierzęta i uprawy. Wejdź na teren za zgodą właściciela.
\end{enumerate}
\end{sekcja}

\begin{sekcja}{Odzyskanie gondoli i danych}
\begin{itemize}
  \item \textbf{Właściciel terenu.} Zanim wejdziesz na teren prywatny, uzyskaj zgodę. Do właściciela idą dwie osoby, rozmowę prowadzi kierownik misji. Krótko wyjaśnij sytuację i pokaż na mapie, gdzie leży gondola. Nie depczcie upraw, zostawcie teren tak, jak był~\cite[rozdz.~2]{nasahb}.
  \item \textbf{Gondola na drzewie.} Nie wspinaj się na drzewo. Wezwij arborystę albo firmę z podnośnikiem. Kontakt do takiej firmy w rejonie lądowania przygotuj dzień przed startem (formularz F8). Po pierwszym locie w 2016 roku gondolę zdjęła z drzewa straż pożarna podnośnikiem (rysunek~\ref{fig:2016drzewo}).
  \item \textbf{Anomalia w locie} (odcinacz nie zadziałał, gondola uszkodzona, brak danych): przed dotknięciem sprzętu zrób zdjęcia ze wszystkich stron. Zapisz, co i jak było połączone. Dopiero potem demontuj~\cite[rozdz.~11]{nasahb}.
  \item \textbf{Dane.} Zaraz po odzyskaniu skopiuj dane z kart pamięci na \textbf{dwa osobne nośniki} (np. laptop i pendrive). Nośniki niech wiozą dwie różne osoby. To nie jest przesada: członkom zespołu zdarzało się zgubić laptop w podróży albo stracić dane z powodu awarii komputera. Po powrocie od razu wgraj dane do repozytorium zespołu. Karty pamięci opisz i montuj tak, aby dało się je wyjąć bez rozbierania gondoli. Odczyt danych z gondoli przećwicz przed misją, w próbie generalnej~\cite[rozdz.~12]{nasahb}.
\end{itemize}
\end{sekcja}
